\begin{table}[!htbp]
\centering
\footnotesize
\caption{Unconditional recodes for familiarity-conditioned Bazooka outcomes}
\label{tab:appendix-unconditional-recodes}
\begin{tabular}{llrrrr}
\hline\hline
Outcome & Coding/sample & N & Trial pack & Cooking demo & Interaction \\
\hline
Recommended Bazooka & Familiar only & 1260 & 44.0 & 11.5 & -11.1 \\
 &  &  & (5.6) & (6.5) & (8.2) \\
 & Unfamiliar coded 0 & 1538 & 51.4 & 11.1 & -8.7 \\
 &  &  & (4.6) & (5.1) & (7.3) \\
 & Lower bound & 1538 & 51.4 & 11.1 & -8.7 \\
 &  &  & (4.6) & (5.1) & (7.3) \\
 & Upper bound & 1538 & 22.1 & 3.0 & -3.3 \\
 &  &  & (4.9) & (5.5) & (7.3) \\
\hline
Perceives Bazooka as risky & Familiar only & 1207 & 2.9 & 0.7 & -4.7 \\
 &  &  & (4.9) & (6.2) & (7.7) \\
 & Unfamiliar coded 0 & 1485 & 6.7 & 1.5 & -5.0 \\
 &  &  & (3.9) & (4.2) & (6.0) \\
 & Lower bound & 1485 & 6.7 & 1.5 & -5.0 \\
 &  &  & (3.9) & (4.2) & (6.0) \\
 & Upper bound & 1485 & -24.5 & -6.6 & 0.5 \\
 &  &  & (6.0) & (7.0) & (8.5) \\
\hline\hline
\end{tabular}
\begin{minipage}{0.92\textwidth}
\scriptsize
Notes: Entries are intent-to-treat effects in percentage points from the fully interacted specification, with CR2 standard errors clustered at the village level in parentheses. ``Familiar only'' reproduces the original skip-pattern sample. For the unconditional recode, respondents unfamiliar with Bazooka are retained and coded as 0. Lower and upper bounds recode unfamiliar respondents as 0 and 1, respectively; ordinary item non-response among familiar respondents is left missing.
\end{minipage}
\end{table}
